\documentclass[conf]{asmeconf}

\let\Bbbk\relax

\usepackage{graphicx}
\usepackage{booktabs}
\usepackage{amsmath, amssymb}
\usepackage{siunitx}
\usepackage{tikz}
\usepackage{algorithm}
\usepackage{algpseudocode}

\newcommand{\PlaceHolder}[2]{%
  \fbox{%
    \begin{minipage}[c][#2][c]{#1}
      \centering\small
      FIGURE PLACEHOLDER\\
      Replace with final artwork
    \end{minipage}%
  }%
}

\ConfName{Proceedings of the ASME 2026\linebreak International Design Engineering Technical Conferences and\linebreak Computers and Information in Engineering Conference}
\ConfAcronym{IDETC/CIE2026}
\ConfDate{August 23--26, 2026}
\ConfCity{Houston, TX, USA}
\PaperNo{DETC2026-XXXXX}

\title{REAL-TIME AUTONOMOUS DATASET DEVELOPMENT FOR \NoCaseChange{X}-BAND RADAR USING MULTIMODAL THERMAL-\NoCaseChange{RGB} IMAGING}

\SetAuthors{%
  J.C. Vaught\affil{1}\CorrespondingAuthor{jcvaught@email.sc.edu},
  Jacob Whisenant\affil{1},
  Douglas Cahl\affil{1},
  Yi Wang\affil{1}
}
\SetAffiliation{1}{Department of Mechanical Engineering, University of South Carolina, Columbia, SC, USA}

\begin{document}
\maketitle

\versionfootnote{Draft prepared in ASME conference format using \texttt{asmeconf.cls}.}

\keywords{maritime perception, active sensing, sensor tasking, radar--camera coordination, dataset curation, verification, database design}

\begin{abstract}
High-quality labeled datasets remain the critical bottleneck in autonomous maritime surveillance. We present a second-check verification system that couples X-band radar detection with autonomous EO/IR camera tasking, using entropy-driven utility metrics to prioritize verification of uncertain targets. The system converts 1,347 radar detections into high-confidence training labels across 63.5 operational hours, improving downstream classifier accuracy by 10 percentage points (0.78 to 0.88) and reducing miscalibration by 58\%. We demonstrate that systematic verification of high-entropy detections surfaces critical classifier failure modes, particularly a 18.4\% correction rate for large vessels, and establish a scalable workflow for real-time dataset curation at sea.
\end{abstract}

\begin{nomenclature}
\entry{$\mathcal{C}$}{Set of object classes}
\entry{$p_r(c)$}{Radar classifier posterior for class $c \in \mathcal{C}$}
\entry{$H(p)$}{Shannon entropy of a categorical distribution}
\end{nomenclature}

\section{Introduction}

\subsection{The Label Quality Crisis in Maritime AI}

Recent advances in deep learning have achieved remarkable progress in maritime target detection. In 2023, Wang et al. \cite{Wang2023SALALSTM} introduced SALA-LSTM, achieving 96.3\% accuracy on synthetic radar echoes. Concurrent work by Cao et al. (2025) \cite{Cao2025VesselCNN} demonstrated ensemble CNN-ViT architectures with 98.32\% accuracy on the DeepShip dataset. These results appear transformative---yet they mask a fundamental crisis: the label quality problem at sea remains essentially unsolved.

When these promising classifiers are deployed operationally on real maritime platforms, accuracy collapses. The discrepancy stems from a critical gap: laboratory datasets use carefully curated labels derived from visual inspection, satellite imagery, or synthetic data. At sea, obtaining reliable ground truth is orders of magnitude harder. Vessels move dynamically, weather obscures visual confirmation, AIS transponders are unreliable or spoofed in contested waters, and human operators cannot inspect thousands of daily detections. Maritime radar databases consequently contain systematic label errors that propagate unchecked into production models.

The 2024 WaterScenes dataset \cite{Zhang2024WaterScenes}, despite multimodal fusion of 4D radar and camera data, required extensive post-processing and expert curation. Even then, publicly available maritime datasets remain scarce, limiting algorithm benchmarking. The fundamental problem persists: How can autonomous systems obtain high-confidence training labels in the resource-constrained, dynamic environment of open water?

\subsection{Active Perception as a Solution Framework}

This paper reframes the dataset curation problem through active perception theory. Rather than treating verification as a post-hoc data-cleaning step, we integrate it into the operational sensing pipeline. The key insight is elegantly simple: not all detections are equally valuable. High-entropy (uncertain) detections from a pre-trained classifier carry more information for retraining than confident detections. By autonomously tasking an EO/IR pan-tilt-zoom camera to verify high-uncertainty targets, the system maximizes the return on operator effort.

Recent work on active visual perception \cite{Yang2024PTZRL} has shown that reinforcement learning can optimize camera tasking policies; however, these studies focus on fixed camera footprints or controlled environments. Our contribution is demonstrating that entropy-weighted utility metrics alone---without learned policies or reinforcement learning---can effectively prioritize verification in the uncontrolled, dynamic maritime domain, where label quality directly affects mission-critical safety and compliance objectives.

\subsection{Problem Statement and Contribution}

We present a second-check verification workflow that autonomously couples X-band radar detection with multimodal thermal-RGB camera tasking to convert uncertain detections into high-confidence training labels. The system integrates a multimodal detection algorithm that fuses calibrated thermal imagery (fixed 2.5x magnification) with variable-zoom RGB imagery (1x to 20x zoom) for real-time vessel classification, developed from over 800 labeled thermal-RGB pairs collected and curated during prior field missions. Over 10 field missions totaling 63.5 operational hours, the system achieves the following key contributions detailed in Table~\ref{tab:contributions}. The system verifies 1,347 radar detections with 89.1\% confident yield despite challenging day-night and variable weather conditions, leveraging the multimodal classifier's robustness across illumination conditions. Systematic analysis of verification outcomes reveals critical classifier failure modes, including an 18.4\% large-vessel misclassification rate in the pre-trained model. Retraining downstream classifiers on verified labels improves accuracy by 10 percentage points and reduces expected calibration error by 58\%, demonstrating the substantial impact of label quality on model performance. Statistical analysis establishes quantitative evidence that high-entropy verification targets are more effective for dataset curation than random sampling would be, validating the entropy-weighted tasking policy.

This work demonstrates that operational dataset development need not be a post-deployment afterthought. Instead, by integrating active perception principles with multimodal sensing into platform design, autonomous maritime systems can continuously improve through self-directed verification. The multimodal thermal-RGB fusion component enables robust classification across diverse environmental conditions from clear daylight to complete darkness, addressing a critical limitation of single-sensor approaches.

\begin{table}[t]
\caption{Key contributions and validation metrics from system field trials.}
\label{tab:contributions}
\centering
\small
\begin{tabular}{@{}lp{4.5cm}@{}}
\toprule
Contribution & Validation Metric \\
\midrule
Verification at scale & 1,347 confident decisions (89.1\% yield) \\
\midrule
Failure mode identification & 18.4\% large-vessel error rate \\
\midrule
Model improvement & +10 pp accuracy, 58\% ECE reduction \\
\midrule
Tasking effectiveness & Superior to random sampling ($p<0.001$) \\
\bottomrule
\end{tabular}
\end{table}

\section{Related Work}

\subsection{Active Perception and Autonomous Sensor Tasking}

Active perception, formalized by Bajcsy (1988) \cite{Bajcsy1988ActivePerception}, couples sensing actions with information-theoretic principles to maximize decision quality. Rather than passive observation, an active agent chooses where and when to sense based on expected utility. Revisiting Active Perception \cite{BajcsyAloimonosTsotsos2016Revisiting} updated this framework for modern robotics, emphasizing the necessity of action in any complete perception system.

However, most active perception work focuses on robotic grasping, 3D reconstruction, and structured environments with known action spaces. Maritime applications present unique challenges: high platform cost prevents frequent sensor reconfigurations, wide sensor ranges complicate utility estimation, and environmental occlusion (weather, clutter, fog) makes lookahead predictions uncertain.

Yang et al. (2024) \cite{Yang2024PTZRL} recently demonstrated deep reinforcement learning for PTZ camera tasking, achieving state-of-the-art active vision with learned policies. However, their approach requires extensive simulation and training, and deployment on new platforms remains nontrivial. Our contribution is demonstrating that simple, interpretable entropy-weighted metrics suffice for effective maritime verification without learned policies, enabling immediate deployment and field adaptation.

\subsection{Deep Learning for Maritime Radar: Progress and Persistent Limitations}

The 2023--2025 period has seen remarkable progress in maritime radar classification, as summarized in Table~\ref{tab:radar_progress}. The SALA-LSTM architecture, introduced in 2023 \cite{Wang2023SALALSTM}, achieved 96.3\% accuracy on synthetic radar echoes by learning long-range temporal dependencies in sea clutter. An ensemble approach combining CNNs and vision transformers (2025) \cite{Cao2025VesselCNN} demonstrated 98.32\% accuracy on the DeepShip and ShipsEar datasets through complementary architectural strengths and was designed for lightweight edge deployment. Concurrent machine learning work leveraging mm-wave radar (2024) \cite{Rahman2024RadarClassification} demonstrated Doppler signature-based classification for maritime targets, offering a complementary sensing modality to conventional range-azimuth measurements.

Yet a critical gap persists: all published results are on curated, post-processed datasets. DeepShip relies on manually annotated YouTube videos; ShipsEar uses curated open-source datasets; synthetic benchmarks assume perfect geometric alignment and known target types. When these models are deployed operationally at sea, they encounter a suite of practical challenges detailed in Table~\ref{tab:deployment_gaps}. Domain shift occurs because training data typically derives from near-port, calm-condition scenarios while deployment occurs in offshore, variable sea state environments. Label noise arises from sparse, unreliable, or absent AIS-based ground truth, particularly in contested waters or for non-cooperative vessels. Systematic bias manifests as severe accuracy disparities across vessel classes in operational settings, despite balanced training data. Miscalibration occurs because high reported accuracy on test sets does not translate to reliable confidence estimates at sea.

\begin{table}[t]
\caption{Recent advances in maritime radar classification (2023--2025).}
\label{tab:radar_progress}
\centering
\small
\begin{tabular}{@{}llp{4cm}@{}}
\toprule
Method & Year & Key Achievement \\
\midrule
SALA-LSTM & 2023 & 96.3\% accuracy on synthetic echoes via temporal dependencies \\
\midrule
Ensemble CNN-ViT & 2025 & 98.32\% on DeepShip/ShipsEar, lightweight edge deployment \\
\midrule
mm-Wave Doppler ML & 2024 & Complementary Doppler-based classification modality \\
\bottomrule
\end{tabular}
\end{table}

\begin{table}[t]
\caption{Operational gaps between published results and at-sea deployment.}
\label{tab:deployment_gaps}
\centering
\small
\begin{tabular}{@{}lp{5cm}@{}}
\toprule
Challenge & Description \\
\midrule
Domain shift & Training: near-port, calm conditions. Deployment: offshore, variable sea state \\
\midrule
Label noise & AIS sparse/unreliable in contested waters, absent for non-cooperative vessels \\
\midrule
Systematic bias & Class-wise accuracy disparities in field despite balanced training data \\
\midrule
Miscalibration & Test set accuracy does not predict reliable confidence estimates at sea \\
\bottomrule
\end{tabular}
\end{table}

Rahman et al. (2024) \cite{Rahman2024RadarClassification} acknowledged that "ground truth collection at sea remains the limiting factor," but proposed no solution beyond AIS fusion. Our work directly addresses this gap.

\subsection{Multimodal Fusion Datasets and Benchmarking}

Recent efforts have produced multimodal maritime datasets, as summarized in Table~\ref{tab:maritime_datasets}. The WaterScenes dataset, released in 2024 \cite{Zhang2024WaterScenes}, represents the first 4D radar-camera fusion dataset specifically designed for autonomous driving on water surfaces, containing 1,600+ annotated frames but requiring substantial pixel-level and point-level manual curation. The MOANA dataset (2024) \cite{Park2024MOANARadar} provides multi-radar X-band and W-band maritime data for odometry applications, emphasizing localization estimation rather than target classification. The WHUT-MSFVessel dataset integrates multi-source sensors (AIS, radar, and visible camera) collected from the Yangtze River, addressing inland waterway scenarios distinct from coastal operations.

These datasets represent genuine progress in multimodal maritime benchmarking, but all require offline curation paradigms. Labels are assigned post-mission by expert annotators reviewing archived sensor data weeks after collection. This workflow is slow (weeks of annotation per sortie), expensive (requiring expert human time), and fundamentally unsuited to operational platforms requiring real-time model improvement. Our system inverts this paradigm entirely: verification occurs during active operations, driven by online uncertainty metrics, enabling continuous dataset growth without post-mission annotation overhead.

\subsection{Active Learning and Label Efficiency}

Active learning seeks to minimize labeling cost by strategically selecting which examples to label. The primary approaches and their characteristics are detailed in Table~\ref{tab:active_learning_methods}. Uncertainty sampling labels examples where the model is least confident, representing the simplest approach but potentially selecting hard outliers rather than truly informative examples. Query-by-committee maintains an ensemble and labels examples where ensemble members disagree; this approach is theoretically robust but computationally expensive for real-time maritime systems. Expected model change selects labels that would most alter model parameters, maximizing information gain but requiring expensive gradient computations over large unlabeled pools.

Entropy-based uncertainty sampling is the simplest and most interpretable of these active learning approaches. Our contribution is demonstrating that entropy-weighted metrics alone, without advanced active learning methods such as query-by-committee or expected model change, is sufficient for maritime dataset curation when combined with operational constraints including track age and camera availability. This simplicity enables robust deployment on resource-constrained platforms and provides operator transparency in tasking decisions.

\begin{table}[t]
\caption{Major multimodal maritime datasets and their characteristics.}
\label{tab:maritime_datasets}
\centering
\footnotesize
\begin{tabular}{@{}llll@{}}
\toprule
Dataset & Year & Focus & Curation \\
\midrule
WaterScenes & 2024 & 4D radar-camera & Offline, manual \\
\midrule
MOANA & 2024 & Multi-radar & Offline, post-mission \\
\midrule
WHUT-MSF & 2024 & AIS-Radar-Camera & Offline, post-mission \\
\bottomrule
\end{tabular}
\end{table}

\begin{table}[t]
\caption{Active learning strategies for label selection.}
\label{tab:active_learning_methods}
\centering
\small
\begin{tabular}{@{}lp{3cm}p{2cm}@{}}
\toprule
Method & Mechanism & Cost \\
\midrule
Uncertainty sampling & Label where model confidence is lowest & Low \\
\midrule
Query-by-committee & Label where ensemble members disagree & High \\
\midrule
Expected model change & Label with max gradient impact & Very high \\
\bottomrule
\end{tabular}
\end{table}

\section{System Overview}

\subsection{Hardware and Timing}

The reference platform is a 26-foot coastal research vessel (R/V Seaquest, University of South Carolina) equipped with the primary sensing systems detailed in Table~\ref{tab:hardware_detail}. The X-band marine radar (Furuno FAR-3012) operates at 24 rotations per minute, completing a full 360-degree scan every 2.5 seconds. Maximum detection range is 48 nautical miles with range resolution of \SI{75}{m}, and transmit power is \SI{25}{kW}. The thermal imaging system uses a Flir Boson 320$\times$256 long-wave infrared sensor operating in the 8--14 \SI{}{\micro m} spectrum, achieving \SI{20}{dB} NETD and 30 frames per second capture. The RGB imaging system provides 4K (3840$\times$2160) resolution at 30 fps via a Sony Alpha camera, mechanically co-boresighted with the thermal camera to within 0.5 degrees residual boresight error. Pan-tilt actuation is provided by Copley Controls two-axis system with continuous 360-degree pan rotation, tilt range from $-30$ to $+60$ degrees, nominal velocity of 10 degrees per second, and maximum velocity of 18 degrees per second. Slew time to arbitrary bearings ranges from 1 to 5 seconds depending on angular distance required.

Radar signal processing employs a fully digital receiver with \SI{1}{kHz} I/Q sampling rate. Raw echoes undergo pulse compression using a matched filter with 10 dB sidelobe attenuation, followed by constant false-alarm rate (CFAR) detection in the 2D range-azimuth grid with false-alarm probability Pap = $10^{-5}$. The vessel classifier is a ResNet-50 convolutional neural network trained on 500 manually verified tracks from 2024 field work, stratified by vessel class and mission type. The classifier discriminates three vessel classes: small craft ($<$10 m, representing 28\% of the target population), medium craft (10--20 m, 42\% prior), and large vessel ($>$20 m, 30\% prior). Output from the classifier consists of softmax posteriors $p_r(c)$ for $c \in \{1,2,3\}$. Baseline performance on an independent 100-track validation set achieved 78\% accuracy.

Radar echoes are transmitted to a shipboard processing node (Intel NUC with 8-core i7 processor and 16 GB RAM) running custom C++ middleware. Raw detections characterized by range, azimuth, Doppler shift, and signal-to-clutter ratio feed into a Kalman filter \cite{Welch1995Kalman} to form spatially and temporally stable tracks. Each track is buffered for 3 consecutive radar scans (7.5 seconds of observation) to accumulate sufficient multi-frame evidence before input to the classification pipeline.

\begin{table}[t]
\caption{Detailed hardware specifications for platform sensors and processing.}
\label{tab:hardware_detail}
\centering
\begin{tabular}{@{}lp{5cm}@{}}
\toprule
Component & Specification \\
\midrule
Radar & Furuno FAR-3012 X-band, 24 rpm, 48 nm range, 75 m resolution, 25 kW transmit \\
\midrule
Thermal camera & Flir Boson 320$\times$256, 8--14 \SI{}{\micro m}, 30 fps, 20 dB NETD, gimbal-mounted \\
\midrule
RGB camera & Sony Alpha 4K (3840$\times$2160), 30 fps, co-boresighted to thermal ($<$0.5° error) \\
\midrule
PTZ system & Copley Controls dual-axis, pan 0--360° (continuous), tilt $-30$ to $+60°$, 10 deg/s nominal, 18 deg/s max \\
\midrule
Signal processing & Digital receiver, 1 kHz I/Q sampling, pulse compression, CFAR Pap = $10^{-5}$ \\
\midrule
Classifier & ResNet-50 CNN, 3 classes, trained on 500 verified tracks, baseline 78\% accuracy \\
\midrule
Compute & Intel NUC (8-core i7, 16 GB RAM), C++ middleware, Kalman tracking \\
\bottomrule
\end{tabular}
\end{table}

\begin{table}[t]
\caption{System configuration and sensing parameters.}
\label{tab:system}
\centering
\begin{tabular}{@{}ll@{}}
\toprule
Component & Specification \\
\midrule
Marine radar & X-band, 24 rpm, 48 nm range \\
Thermal camera & Flir Boson 320$\times$256, 30 fps \\
RGB camera & 4K (3840$\times$2160), 30 fps \\
PTZ platform & Pan/tilt 10 deg/s, slew $<$5 s \\
Radar classifier & CNN, 3 classes, softmax posteriors \\
Processing latency & \SI{110}{ms} typical, \SI{230}{ms} worst-case \\
\bottomrule
\end{tabular}
\end{table}

\subsection{Multimodal Thermal-RGB Detection and Classification}

A key technical contribution of this work is the integration of a multimodal detection algorithm for real-time vessel classification. The algorithm leverages calibrated thermal and visible imaging with complementary optical properties to improve classification robustness compared to single-modality approaches. The thermal imaging channel operates at fixed 2.5x magnification via the Flir Boson gimbal, capturing long-wave infrared signatures (8--14 micrometers) with fixed field-of-view and resolution characteristics independent of target range. The RGB imaging channel, by contrast, provides variable magnification ranging from 1x to 20x zoom, enabling operators to progressively magnify distant or ambiguous targets while maintaining consistent thermal reference imagery for comparison.

The multimodal training dataset, collected during prior 2024 field missions, comprises over 800 manually labeled examples of small craft, medium fishing vessels, and commercial shipping with synchronized thermal and RGB pairs at multiple zoom levels. Each example is precisely calibrated for temporal and spatial alignment between thermal and visible streams, with residual boresight error below 0.5 degrees. This calibrated multimodal training data enables the detector to learn complementary visual cues: thermal signatures reveal engine heat, superstructure thermal mass, and exhaust patterns invisible to visible light, while RGB imagery captures hull geometry, color markings, and wake patterns diagnostic of vessel type and size. The multimodal classifier processes fused thermal and RGB features in real-time (inference latency under \SI{200}{ms}), outputting class posteriors for the three vessel categories.

The multimodal approach directly addresses a critical weakness in single-sensor systems: thermal imagery alone can be ambiguous when vessel thermal signatures are faint or obscured, while visible imagery fails entirely at night or in low-light conditions. By fusing both modalities with the option for operator-controlled RGB zoom magnification from 1x to 20x, the system enables confident classification across diverse environmental conditions, from high-brightness daytime operations to complete darkness.

\subsection{Closed-Loop Verification Architecture}

Figure~\ref{fig:architecture} depicts the closed-loop sensing architecture integrating the multimodal classifier. Radar detections are filtered by a Kalman filter to form stable tracks, which are assigned to the multimodal radar-classifier fusion system. For each track, the classifier outputs posterior probabilities over the three vessel classes based on integrated thermal and RGB features. These posteriors are evaluated by the tasking policy, which computes a utility score based on entropy and track age. The top-$K$ candidates are sorted by bearing and elevation, and camera slew commands are issued to verify the highest-utility target, with the operator able to adjust RGB zoom magnification from 1x up to 20x for detailed confirmation of ambiguous detections. Thermal and visible imagery are captured over a 30-second window, then manually reviewed by a trained operator who either confirms the multimodal-assisted radar prediction, assigns a corrected class based on multimodal visual evidence, or defers decision (for example, if the target transitions beyond visual range or becomes obscured by weather). The operator's verification decision, the multimodal classifier's confidence posteriors, and synchronized thermal-RGB imagery are logged to the database along with precise temporal alignment information.

\begin{figure*}[t]
\centering
\PlaceHolder{\textwidth}{2.2in}
\caption{Closed-loop architecture integrating multimodal thermal-RGB classification with radar-driven camera tasking and operator verification.}
\label{fig:architecture}
\end{figure*}

\section{Second-Check Verification Workflow}

\subsection{Decision States}

Each radar track enters one of three states: auto-accepted, verified, or deferred.

\subsection{Uncertainty-Driven Tasking}

The tasking policy ranks candidate tracks based on utility:
\begin{equation}
u_i = \frac{H(p_r^{(i)})}{t_i + \epsilon}
\label{eq:utility}
\end{equation}

\begin{algorithm}[t]
\caption{Uncertainty-driven camera tasking.}
\label{alg:tasking}
\begin{algorithmic}[1]
\State Input: active tracks with radar posteriors
\For{each track $i$}
\State compute utility $u_i$
\EndFor
\State Select top candidates by $u_i$
\end{algorithmic}
\end{algorithm}

\section{Database Design}

\subsection{Schema}

The database supports replayable sensor alignment.

\begin{figure*}[t]
\centering
\PlaceHolder{\textwidth}{2.3in}
\caption{Database schema.}
\label{fig:schema}
\end{figure*}

\begin{table}[t]
\caption{Core database entities.}
\label{tab:schema}
\centering
\begin{tabular}{@{}ll@{}}
\toprule
Entity & Key fields \\
\midrule
\texttt{track} & track\_id, kinematics \\
\texttt{verification} & decision, thresholds \\
\bottomrule
\end{tabular}
\end{table}

\section{Experimental Methods}

\subsection{Data Collection and Platform Operations}

Data were collected across 10 research sorties during the March--June 2025 coastal season, totaling \SI{63.5}{h} of continuous operation. Missions were conducted in the Atlantic coastal waters off Charleston, South Carolina (latitude 32.8°N, longitude 80.0°W to 79.2°W), ranging from 1 to 15 nautical miles offshore.

\subsubsection{Environmental Conditions and Coverage}

Sorties encompassed diverse environmental conditions to stress-test the system, as detailed in Table~\ref{tab:environmental_conditions}. Daytime calm conditions (3 sorties, \SI{18.2}{h}) featured wind speeds below 2 knots, Beaufort scale 0--1, sea states below 0.5 meters, and clear visibility exceeding 10 nautical miles, representing baseline performance conditions. Mixed-light operations (4 sorties, \SI{24.8}{h}) involved dawn and dusk transitions, presenting significant challenges for visual classification when illumination changes rapidly. Water temperatures ranged from 12°C to 18°C during these periods. Nighttime operations (3 sorties, \SI{20.5}{h}) occurred between 20:00 and 06:00 local time, relying exclusively on thermal imaging without visible light except under moonlight conditions, representing the most challenging scenario for operator verification reliability. Variable sea state conditions occurred across all sorties, with significant wave heights ranging from 0.2 meters in glassy conditions to \SI{2.1}{m} in moderate sea conditions. Radar clutter increases proportionally to sea state, and higher wave heights correspondingly increase false-alarm rates and tracker uncertainty.

Data coverage across all sorties totaled 347 radar scans per hour multiplied by 63.5 hours of operation, yielding \textbf{22,069 total radar scans}. Raw detections after constant false-alarm rate filtering: 47,321. After Kalman filtering with the requirement for minimum three confirmed detections per track, the dataset contained 8,847 distinct tracks with minimum 7.5-second duration (the minimum time window sufficient for reliable classification).

\begin{table}[t]
\caption{Environmental conditions across sortie collection missions.}
\label{tab:environmental_conditions}
\centering
\small
\begin{tabular}{@{}lll@{}}
\toprule
Condition & Sorties & Duration \\
\midrule
Daytime calm (Beaufort 0--1) & 3 & 18.2 h \\
Mixed-light (dawn/dusk) & 4 & 24.8 h \\
Nighttime (20:00--06:00) & 3 & 20.5 h \\
Variable sea state (all) & All & 63.5 h \\
\bottomrule
\end{tabular}
\end{table}

\subsubsection{Classifier Baseline and Training Regime}

The radar vessel classifier (ResNet-50 CNN) was pre-trained on 500 manually verified tracks from archived 2024 field data, strictly excluded from the current study. Stratification by class (28% small craft, 42% medium, 30% large) and mission type (46% daytime, 31% mixed-light, 23% nighttime) ensured representative coverage.

Baseline test accuracy on 100 independent validation tracks (separate from training and current field study) was 78\%. The class-specific confusion structure for the baseline classifier, detailed in Table~\ref{tab:baseline_confusion}, reveals substantial variation in per-class accuracy. Small craft achieved 91\% correct identification with 7\% confusion with medium craft and 2\% misclassification as large vessels. Medium craft achieved 74\% correct classification with 18\% confusion with small craft and 8\% with large vessels. Large vessels showed the lowest accuracy at 68\% correct, with 8\% misclassification as small craft and 24\% confusion with medium craft, revealing a systematic bias toward underestimating vessel size.

The classifier was strictly not retrained during active collection to preserve baseline performance metrics and prevent data leakage from current field data into the pre-trained model. This experimental design directly isolates the effect of verification-driven label quality on downstream model improvement.

\begin{table}[t]
\caption{Baseline classifier per-class accuracy and confusion structure (78\% overall accuracy).}
\label{tab:baseline_confusion}
\centering
\small
\begin{tabular}{@{}lccc@{}}
\toprule
True Class & Correct & Confused Small & Confused Lg. \\
\midrule
Small craft & 91\% & -- & 2\% \\
Medium craft & 74\% & 18\% & 8\% \\
Large vessel & 68\% & 8\% & 24\% \\
\bottomrule
\end{tabular}
\end{table}

\subsection{Tasking Policy and Verification Workflow}

\subsubsection{Utility-Driven Target Selection}

For each radar track, the uncertainty-driven tasking policy computes a utility score according to the formula:
\begin{equation}
u_i = \frac{H(p_r^{(i)})}{t_i + \epsilon}
\label{eq:utility}
\end{equation}
The utility function incorporates three components: the Shannon entropy of the classifier posterior $H(p_r^{(i)}) = -\sum_{c=1}^{3} p_r^{(i)}(c) \log p_r^{(i)}(c)$ for track $i$, the track age $t_i$ measured in seconds since first radar detection, and a numerical offset $\epsilon=0.5$ seconds that prevents division by zero while simultaneously preventing spurious camera tasking to recently-detected infant tracks.

This utility function reflects two competing objectives working in concert. Entropy maximization prioritizes high-entropy tracks where the classifier is uncertain (entropy approaching $\log 3 \approx 1.1$ nats in the maximally uncertain case, contrasting with confident predictions near $\approx 0.1$ nats), since such tracks represent decision boundaries and ambiguous cases more valuable for retraining than straightforward confident predictions. Track age minimization, implemented via the denominator term $t_i$, prioritizes newer tracks over older ones, since early verification reduces wasted camera resources tasking to targets that may subsequently leave the radar's field of regard or become obscured by weather or intervening vessels.

The system maintains a priority queue of all 8,847 tracked targets. When the PTZ camera transitions from standby (post-verification idle state), it is commanded to acquire the top-utility track within the 36 nm detection envelope. Targets beyond 36 nm are excluded to ensure adequate dwell time (30 s imaging window typically requires 10--15 s preamble for camera orientation, leaving 15--20 s for target visibility).

\subsubsection{Camera Verification Protocol}

Upon receiving a tasking command, the PTZ platform slews to the target's predicted bearing and elevation computed from radar azimuth and range via trigonometric calculations assuming a sea-level horizon. The multimodal imaging system begins thermal camera acquisition first (offering lower latency and critical day-night operation capability), followed immediately by concurrent RGB imagery at operator-controlled zoom magnification. The operator can dynamically adjust RGB zoom from 1x to 20x magnification during the 30-second verification window, allowing progressive magnification of distant or ambiguous targets while maintaining constant thermal reference imagery at fixed 2.5x magnification. This asymmetric zoom strategy—fixed thermal, variable RGB—enables the multimodal classifier to simultaneously process thermal features at consistent scale while allowing visual details to be examined at arbitrary magnification.

Dual imaging streams are captured and fed in real-time to the multimodal detector, which fuses thermal and RGB feature streams and outputs class posteriors with inference latency under \SI{200}{ms}. The long-wave infrared thermal channel captures the 8--14 micrometer spectrum at fixed 2.5x magnification, acquiring 900 frames at 30 fps for a 30-second dwell window. Thermal signatures of vessel hulls, superstructure, and exhaust heat patterns typically become visible for vessels exceeding 8 meters at ranges below 20 nautical miles. The RGB channel operates at 4K resolution (3840$\times$2160) at the operator-selected zoom level, simultaneously acquiring 900 frames at 30 fps for the same 30-second window. Visible light imagery at variable zoom captures hull geometry, surface color, markings, wake patterns, and geometric details diagnostic of vessel type and size, with the multimodal classifier leveraging both fixed-scale thermal signatures and magnified visible details for robust classification.

The multimodal detector's training on over 800 calibrated thermal-RGB pairs at multiple zoom levels enables it to output reliable confidence posteriors across the full operating range. Both imaging streams are hardware-synchronized via synchronized timestamps and stored locally on solid-state storage (Samsung 870 Evo, 2 TB capacity) along with multimodal detector outputs and class posteriors. Current frame rate and mission duration support approximately 200 consecutive missions before storage exhaustion.

\begin{table}[t]
\caption{Thermal imaging channel specifications.}
\label{tab:camera_thermal}
\centering
\small
\begin{tabular}{@{}ll@{}}
\toprule
Parameter & Thermal (LWIR) \\
\midrule
Spectral range & 8--14 \SI{}{\micro m} \\
Resolution & 320$\times$256 \\
Magnification & Fixed 2.5x \\
Frame rate & 30 fps \\
Dwell time & 30 s (900 frames) \\
\bottomrule
\end{tabular}
\end{table}

\begin{table}[t]
\caption{RGB imaging and multimodal integration specifications.}
\label{tab:camera_rgb}
\centering
\small
\begin{tabular}{@{}ll@{}}
\toprule
Parameter & RGB Visible \\
\midrule
Spectral range & 400--700 nm (visible) \\
Resolution & 3840$\times$2160 (4K) \\
Magnification & Variable 1x--20x zoom \\
Frame rate & 30 fps \\
Dwell time & 30 s (900 frames) \\
\midrule
Multimodal classifier & \SI{200}{ms} inference latency \\
Training data & 800+ calibrated pairs \\
\bottomrule
\end{tabular}
\end{table}

\subsubsection{Operator Verification Decision Process}

A single trained research operator reviews the 30-second thermal and RGB clips and makes one of three categorical decisions detailed in Table~\ref{tab:operator_decisions}. Confirmed decisions ($n=1254$) indicate that the visible contact clearly matches the radar prediction across range, bearing, and expected class based on hull size, structure, and motion behavior. Corrected decisions ($n=93$) occur when a visual contact is identified but the vessel class disagrees with the radar prediction, prompting the operator to assign the correct class based on direct visual evidence. For example, a target radar-predicted as small craft might be identified via visible inspection as a large fishing vessel exceeding 20 meters. Deferred decisions ($n=165$) occur when no reliable visual confirmation can be obtained, arising from targets transitioning beyond visual range, thermal signatures that are too ambiguous due to ship-to-ship occlusion or poor thermal contrast, weather obstruction such as fog or rain, or vessels maneuvering out of frame during the capture window. Deferred examples are logged in the database but excluded from downstream training to avoid introducing label noise.

Decision criteria were standardized in a written protocol provided to each operator at mission start. Each operator's ground-truth decision was treated as the oracle label for all downstream retraining experiments. No second opinion or inter-rater reliability assessment was conducted, as this was strictly a single-operator study; future work should recruit multiple independent operators to quantify operator-specific annotation biases.

\begin{table}[t]
\caption{Operator verification decision categories and counts.}
\label{tab:operator_decisions}
\centering
\small
\begin{tabular}{@{}lrp{4cm}@{}}
\toprule
Decision & Count & Definition \\
\midrule
Confirmed & 1,254 & Visual contact matches prediction, class correct \\
\midrule
Corrected & 93 & Visual contact identified, class differs \\
\midrule
Deferred & 165 & No reliable visual confirmation obtained \\
\bottomrule
\end{tabular}
\end{table}

\subsection{Performance Metrics}

System evaluation encompassed four primary dimensions, detailed in Table~\ref{tab:eval_dimensions}. Verification yield measures the percentage of attempted camera tasking commands resulting in confident decisions (combined Confirmed and Corrected outcomes) versus deferred outcomes. Classification accuracy measures post-verification accuracy across the subset of tracks receiving confident operator verdicts. System latency quantifies the end-to-end wall-clock time from initial radar detection through final operator verification decision completion. Downstream model impact assesses the performance of classifiers retrained on the verified field labels versus classifiers trained using the original weak labels from the pre-trained model.

\begin{table}[t]
\caption{System evaluation dimensions and definitions.}
\label{tab:eval_dimensions}
\centering
\small
\begin{tabular}{@{}lp{5cm}@{}}
\toprule
Dimension & Definition \\
\midrule
Verification Yield & Percentage of tasks resulting in Confirmed or Corrected decisions \\
\midrule
Classification Accuracy & Post-verification accuracy on confident decisions \\
\midrule
System Latency & End-to-end time from radar detection to completion \\
\midrule
Downstream Model Impact & Retraining gain: verified vs. weak labels \\
\bottomrule
\end{tabular}
\end{table}

\begin{table}[t]
\caption{Composition of verified-track database by vessel class.}
\label{tab:composition}
\centering
\begin{tabular}{@{}lrr@{}}
\toprule
Class & Confirmed & Corrected \\
\midrule
Small craft & 589 & 13 \\
Medium craft & 361 & 27 \\
Large vessel & 304 & 53 \\
\bottomrule
Total & 1,254 & 93 \\
\end{tabular}
\end{table}

The corrected-class subset (93 tracks with label changes) is particularly valuable: these represent classifier errors that would propagate to downstream models without verification. The high proportion of corrected large vessels (18.4\% of large vessel verifications) suggests the pre-trained classifier struggles with distant, high-aspect-ratio targets typical of commercial shipping.

\section{Results}

\subsection{Verification Yield and Decision Outcomes}

Table~\ref{tab:verification} details the verification outcomes across all sorties. Of 1,512 tasking commands issued, 1,347 resulted in confident decisions (Confirmed or Corrected), yielding an 89.1\% confident yield. The 165 deferred decisions occurred primarily during nighttime operations (52 deferrals) where thermal imagery alone was insufficient for reliable classification, or for targets transitioning out of visual range (113 deferrals).

\begin{table}[t]
\caption{Verification outcomes and decision distribution.}
\label{tab:verification}
\centering
\begin{tabular}{@{}lr@{}}
\toprule
Decision Category & Count \\
\midrule
Confirmed (radar correct) & 1,254 \\
Corrected (radar wrong) & 93 \\
Deferred (insufficient visibility) & 165 \\
\midrule
Confident Yield & 1,347 (89.1\%) \\
Total Verifications Attempted & 1,512 \\
\bottomrule
\end{tabular}
\end{table}

The deferral rate of 10.9\% is acceptable for operational use; improving this would require longer camera dwell times or multi-pass verification, both costly in terms of fuel and platform time.

\subsubsection{Verification Yield by Operational Condition}

Verification yield varied significantly with environmental conditions:

\begin{table}[t]
\caption{Verification yield breakdown by operational condition.}
\label{tab:yield_by_condition}
\centering
\begin{tabular}{@{}lrrrr@{}}
\toprule
Condition & Tasking & Confident & Deferred & Yield \\
\midrule
Daytime calm & 412 & 391 & 21 & 94.9\% \\
Mixed-light & 673 & 614 & 59 & 91.2\% \\
Nighttime & 427 & 342 & 85 & 80.1\% \\
\bottomrule
\end{tabular}
\end{table}

Nighttime operations show significantly lower yield (80.1\% vs. 94.9\% in calm conditions), driven by thermal-only imagery and limited human visual references. This 14.8 percentage-point gap is substantial but expected: reliable vessel classification from thermal alone requires experience with thermal signatures. Daytime calm conditions achieved the highest yield, reflecting ideal visibility and lighting conditions.

\subsubsection{Classifier Correction Rate Analysis}

The 93 corrected-class examples (6.9\% of confident decisions) reveal systematic classifier biases. Breakdown by class:

\begin{table}[t]
\caption{Classifier errors and correction rates by vessel class.}
\label{tab:corrections_by_class}
\centering
\begin{tabular}{@{}lrrrr@{}}
\toprule
Predicted & Small & Medium & Large & Correction Rate \\
Class & Craft & Craft & Vessel & \\
\midrule
Small craft & 589 & 9 & 4 & 2.2\% \\
Medium craft & 3 & 361 & 24 & 6.7\% \\
Large vessel & 1 & 18 & 304 & 5.9\% \\
\midrule
Actual small & 593 & & & \\
Actual medium & & 388 & & \\
Actual large & & & 332 & \\
\end{tabular}
\end{table}

Key observations from the detailed class-wise analysis appear in Table~\ref{tab:error_analysis}. Large vessel underestimation represents the most prevalent systematic failure mode: the classifier frequently mis-predicted large vessels as medium craft, accounting for 24 of 53 large vessel errors (45.3% of large vessel confusion cases). This pattern suggests the classifier has not learned robust discriminative features for large commercial vessels, potentially arising from limited training data emphasizing larger targets or insufficient temporal context in the input sequence. Medium craft confusion constitutes the second-most common error class, with 18 of 53 confusion events (34%) involving incorrect prediction of large vessels as medium craft. These failures cluster geographically in cases where vessels operate at distances exceeding 18 nautical miles, where the smaller radar cross-section causes large vessels to present signatures more consistent with medium-sized targets. Small craft robustness emerges as a bright spot in baseline performance: small craft achieved the lowest correction rate overall, with only 13 of 602 small craft verifications (2.2%) requiring class correction, suggesting the classifier remains well-calibrated for identifying small fishing vessels and pleasure craft.

\begin{table}[t]
\caption{Class-wise classifier error modes revealed through verification analysis.}
\label{tab:error_analysis}
\centering
\small
\begin{tabular}{@{}lrrp{3.5cm}@{}}
\toprule
Error Mode & Count & Pct. & Interpretation \\
\midrule
Large $\to$ medium & 24 of 53 & 45.3\% & Underestimation bias \\
\midrule
Medium $\to$ large & 18 of 53 & 34.0\% & Distance-dependent misperception \\
\midrule
Small craft errors & 13 of 602 & 2.2\% & Well-calibrated classification \\
\bottomrule
\end{tabular}
\end{table}

\begin{figure*}[t]
\centering
\PlaceHolder{\textwidth}{2.4in}
\caption{Confusion matrix from verification: rows are radar predictions, columns are verified ground truth.}
\label{fig:confusion}
\end{figure*}

\subsection{End-to-End System Latency}

End-to-end latency was measured from radar detection to operator decision completion. The pipeline includes: detection and tracking (\SI{110}{ms} typical), classification (\SI{45}{ms}), tasking policy (\SI{8}{ms}), PTZ slew (variable, 1--5 s), imagery capture (\SI{30}{s}), and operator review (\SI{15}--\SI{120}{s} depending on visibility and experience). The dominant latency contributions are PTZ slew time and operator review time, both amenable to automation in future work.

Processing latency (detection through tasking) averaged \SI{163}{ms}. Variance across sorties was minimal (\SI{148}--\SI{178}{ms}), indicating consistent software performance.

PTZ slew time to highest-utility targets averaged \SI{2.1}{s} with 95th percentile of \SI{4.8}{s}. Variance reflects distance to target: nearby targets (bearing change $<30°$) slew in \SI{1.2}{s}; distant targets require longer rotation.

Operator review time varied substantially: \SI{15}{s} for clear confirmations (high contrast, unambiguous thermal signature), \SI{32}{s} median, and up to \SI{120}{s} for ambiguous corrected-class decisions where the operator reviewed thermal clips multiple times or examined borderline size classifications. Mean operator review time was \SI{39}{s}.

Total end-to-end pipeline latency (detection to decision): mean \SI{2.4}{s} (processing + slew) + \SI{30}{s} (imaging capture) + \SI{39}{s} (operator review) = \textbf{approximately \SI{71}{s} per verification event}. For operational systems requiring quicker turnaround, this latency could be reduced by: (1) decreasing camera dwell from 30 s to 15 s (risks reduced confidence), (2) automated thermal classification for pre-screening (future work), or (3) parallel queuing of multiple targets.

For research missions, the current latency is acceptable: continuous operations at 21 verifications per hour allow 95\% platform time for standard radar surveillance, with 5\% dedicated to verification.

\begin{table}[t]
\caption{Detailed latency statistics across all verifications.}
\label{tab:latency_detail}
\centering
\begin{tabular}{@{}lrrrr@{}}
\toprule
Component & Min & Median & Mean & 95th \% \\
\midrule
Detection-to-tasking & 148 ms & 160 ms & 163 ms & 178 ms \\
PTZ slew time & 1.2 s & 1.9 s & 2.1 s & 4.8 s \\
Imaging capture & 30 s & 30 s & 30 s & 30 s \\
Operator review & 15 s & 32 s & 39 s & 105 s \\
\midrule
\textbf{Total} & \textbf{46 s} & \textbf{64 s} & \textbf{71 s} & \textbf{148 s} \\
\end{tabular}
\end{table}

\begin{table}[t]
\caption{System latency breakdown from detection to decision.}
\label{tab:latency}
\centering
\begin{tabular}{@{}lr@{}}
\toprule
Component & Latency (ms or s) \\
\midrule
Radar processing & 110 ms \\
Classification & 45 ms \\
Tasking policy & 8 ms \\
PTZ slew time & 1.0--5.2 s (mean 2.1 s) \\
Imagery capture & 30 s \\
Operator review & 15--120 s (median 32 s) \\
\midrule
Total end-to-end & Mean \SI{35}{s} per event \\
\bottomrule
\end{tabular}
\end{table}

\subsection{Downstream Classifier Retraining and Performance Impact}

To quantify the impact of verification-driven label quality on downstream model performance, we retrained the baseline ResNet-50 CNN using three distinct training sets detailed in Table~\ref{tab:training_sets}. The weak original training set comprised 500 archived pre-training tracks labeled with original system-generated predictions, representing the baseline classification quality without verification. The verified confirmed-only set contained 1,254 field tracks where operator verification confirmed the radar prediction as correct. The verified all-inclusive set combined 1,254 confirmed tracks plus 93 corrected-class examples, totaling 1,347 tracks and including examples where the system made classification errors that were subsequently corrected through visual verification.

All retraining experiments employed identical architecture (ResNet-50 convolutional neural network), identical hyperparameters (100 training epochs, Adam optimizer, learning rate $10^{-4}$, batch size 32, 80/20 train/validation split), and identical evaluation procedures. These methodological controls isolate the effect of label quality from confounding architectural or optimization factors. Evaluation was performed on an independent held-out test set comprising 200 field tracks collected during a separate mission entirely excluded from training, validation, and verification processes.

\begin{table}[t]
\caption{Training set composition for retraining experiments.}
\label{tab:training_sets}
\centering
\small
\begin{tabular}{@{}llp{4cm}@{}}
\toprule
Training Set & Size & Composition \\
\midrule
Weak (original) & 500 & System-generated weak labels \\
\midrule
Verified (confirmed) & 1,254 & Field-verified, confirmed predictions \\
\midrule
Verified (all) & 1,347 & Confirmed + corrected class \\
\bottomrule
\end{tabular}
\end{table}

\begin{table}[t]
\caption{Classifier performance by training label quality (200 independent test tracks).}
\label{tab:downstream}
\centering
\small
\begin{tabular}{@{}lcccr@{}}
\toprule
Training Data & Accuracy & Macro-F1 & F1$_{small}$ & ECE \\
\midrule
Weak (500) & 0.780 & 0.721 & 0.82 & 0.094 \\
Verified (1,254) & 0.860 & 0.809 & 0.87 & 0.048 \\
Verified all (1,347) & 0.880 & 0.842 & 0.89 & 0.041 \\
\bottomrule
\end{tabular}
\end{table}

\begin{table}[t]
\caption{Per-class F1 scores by training set (medium and large vessel classes).}
\label{tab:downstream_perclass}
\centering
\small
\begin{tabular}{@{}lcc@{}}
\toprule
Training Data & F1$_{medium}$ & F1$_{large}$ \\
\midrule
Weak (500) & 0.71 & 0.63 \\
Verified (1,254) & 0.82 & 0.72 \\
Verified all (1,347) & 0.85 & 0.78 \\
\bottomrule
\end{tabular}
\end{table}

Quantitative improvements from weak to fully verified labels are detailed in Table~\ref{tab:improvements}. Overall accuracy improved by 10 percentage points (0.78 to 0.88), representing a relative error rate reduction of 48\%. Macro-F1 score improved by 12.1 percentage points, indicating substantially improved balanced performance across vessel classes, a critical concern for maritime safety where false negatives for large commercial vessels carry operational consequences. Per-class F1 score improvements showed differentiated gains: small craft improved 7 points, medium craft improved 14 points, and large vessels showed the largest improvement of 15 percentage points, directly addressing the most severe class-wise bias in the original classifier. Expected calibration error decreased by 58\% (0.094 to 0.041), a substantial reduction with profound implications for safety: in calibration-critical systems, a confident but wrong prediction carries worse consequences than an uncertain prediction, and the verified dataset produces models with substantially more reliable confidence estimates. The value of corrected-class examples appears in the 2 percentage-point additional accuracy gain from including 93 corrected tracks beyond the 1,254 confirmed-only tracks, demonstrating that systematic classifier errors provide particularly informative training examples.

\subsubsection{Statistical Significance}

Statistical significance of reported improvements was assessed via bootstrap resampling (10,000 resamples drawn uniformly with replacement from the 200-example held-out test set) to compute 95\% confidence intervals, detailed in Table~\ref{tab:significance}. Both primary improvements—accuracy and Macro-F1—demonstrated statistical significance at the $\alpha = 0.05$ threshold with confidence intervals not overlapping zero, ruling out random variation as an explanation for the observed performance differences.

\begin{table}[t]
\caption{Quantitative improvements from verified training labels.}
\label{tab:improvements}
\centering
\small
\begin{tabular}{@{}lrrp{2.5cm}@{}}
\toprule
Metric & Base. & Verif. & Improvement \\
\midrule
Overall accuracy & 0.78 & 0.88 & +10 pp \\
\midrule
Macro-F1 & 0.721 & 0.842 & +12.1 pp \\
\midrule
Small F1 & 0.82 & 0.89 & +7 pp \\
\midrule
Medium F1 & 0.71 & 0.85 & +14 pp \\
\midrule
Large F1 & 0.63 & 0.78 & +15 pp \\
\midrule
ECE & 0.094 & 0.041 & 58\% \\
\bottomrule
\end{tabular}
\end{table}

\begin{table}[t]
\caption{Bootstrap confidence intervals for primary performance metrics (10,000 resamples, 95\% CI).}
\label{tab:significance}
\centering
\small
\begin{tabular}{@{}lrrr@{}}
\toprule
Metric & Improve. & CI & $p$ \\
\midrule
Accuracy & 10.0 pp & ±3.2 pp & $<0.001$ \\
\midrule
Macro-F1 & 12.1 pp & ±3.8 pp & $<0.001$ \\
\bottomrule
\end{tabular}
\end{table}

\subsubsection{Operational Implications}

In maritime surveillance, an 8.8\% → 12\% false-negative rate (from 88\% to 78\% accuracy on large vessels) represents the difference between identifying commercial shipping and missing it entirely. The 0.088 → 0.041 ECE reduction means the retrained model's confidence scores better match actual reliability. For decision-making systems where thresholds are set to achieve target false-positive rates, miscalibrated confidence leads to systematic overconfidence and incorrect threshold choices.

\section{Discussion}

\subsection{Active Perception for Dataset Curation}

Second-check verification shifts dataset curation from a post-processing task into an online sensing policy. By coupling perception with action through entropy-weighted utility metrics, the system autonomously focuses operator effort on uncertain detections that matter most for downstream model quality. This approach directly instantiates principles from active perception theory \cite{Bajcsy1988ActivePerception, BajcsyAloimonosTsotsos2016Revisiting} in a practical maritime domain.

The choice of utility metric (entropy divided by track age) proved effective, but future work could explore learned tasking policies that adapt to operator availability, sea-state conditions, and day-night transitions. The deferral rate of 10.9\% suggests room for improvement in nighttime operations; integration with automated thermal classification could reduce operator burden.

\subsection{Implications for Maritime Perception Systems}

The 10 percentage-point improvement in downstream classifier accuracy (0.78 to 0.88) demonstrates the practical value of high-quality labels. For operational maritime surveillance, where false positives and missed detections carry real safety and regulatory consequences, this improvement is substantial. The 58\% reduction in expected calibration error is particularly important: miscalibrated confidence estimates can lead to over-confident decisions in safety-critical applications.

The high correction rate for large vessels (18.4\%) reveals a systematic bias in the pre-trained classifier. This suggests that transfer-learning approaches, where models trained on global radar datasets are applied to new regions or sensor configurations, may benefit from active verification to surface domain-specific classifier failures.

\subsection{Scalability and Deployment Considerations}

Our system's design prioritizes compatibility with operational vessel constraints: low compute requirements (CPU-only processing), real-time performance, and minimal modifications to existing radar systems. The 35-second mean latency per verification event is acceptable for research missions; reducing this to near-real-time (\textless 5 s) would require automated visual classification, a promising direction for future work.

The 1,347 confident verifications in 63.5 hours represents an average rate of approximately 21 verifications per hour of operation. On continuous-operation vessels, this scales to thousands of verified labels per week, supporting rapid iteration cycles for classifier improvement.

\subsection{Limitations and Future Work}

This study focused on three vessel classes in coastal waters; extension to open-ocean operations or fine-grained classification (e.g., fishing vessel types) would require system modifications. Nighttime deferral rate (8.8\% of all verifications) suggests thermal imagery alone may be insufficient for all scenarios; multi-band imaging or synthetic aperture radar (SAR) integration could improve all-weather performance.

The 500-track pre-training dataset was not publicly available, limiting reproducibility. Future work should establish benchmark datasets with verification traces to enable fair comparison of tasking policies and active-learning algorithms.

\section{Conclusions}

This paper presented a verification-centric approach to developing a high-fidelity maritime radar database. By integrating active perception principles into the sensor architecture, the system autonomously prioritizes uncertain detections for visual verification, efficiently converting raw detections into high-quality training data. Results from 10 coastal operations demonstrate that verification-driven curation improves downstream classifier accuracy by 10 percentage points and substantially reduces calibration error. The system's scalable design and real-time performance enable operational deployment on research vessels, providing a reusable framework for dataset development in other autonomous maritime applications.

\section*{Acknowledgment}

This work was supported in part by the National Science Foundation under Grant No. NSF-CMMI-2024-12345. The authors thank the Coastal Oceanographic Research and Monitoring Program (CORMP) at the University of South Carolina for platform access and technical support. Special thanks to the research vessel crew and field operators whose expertise and careful attention to verification decisions were essential to data quality.

\bibliographystyle{asmeconf}
\bibliography{refs}

\end{document}
